\documentclass[11pt]{article}
\usepackage[margin=1in]{geometry}
\usepackage{amsmath,amssymb,amsthm}
\usepackage{algorithm}
\usepackage{algpseudocode}
\usepackage{booktabs}
\usepackage{hyperref}

\newtheorem{theorem}{Theorem}
% Upright body text for definitions/examples/remarks: the worked example spans
% several paragraphs and is unreadable in the default italic theorem style.
\theoremstyle{definition}
\newtheorem{definition}{Definition}
\newtheorem{example}{Example}
\newtheorem{remark}{Remark}

% Local copies of the shared macros in docs/macros.tex (same definitions), so
% this document stays standalone while text can be moved verbatim to/from
% docs/neurips_2026.tex.
\newcommand{\cK}{\mathcal{K}}
\newcommand{\bX}{\mathbf{X}}
\newcommand{\bY}{\mathbf{Y}}
\newcommand{\bZ}{\mathbf{Z}}
\newcommand{\by}{\mathbf{y}}
\newcommand{\bz}{\mathbf{z}}
\newcommand{\bo}{\mathbf{o}}
\newcommand{\pa}{\mathrm{pa}}
\newcommand{\scope}{\mathrm{scope}}
\newcommand{\ext}{\mathrm{ext}}
\newcommand{\prune}{\mathrm{prune}}

\title{Credal Variable Elimination:\\Exact Marginal Inference for Credal Networks}
\author{IBM Research}
\date{}

\begin{document}
\maketitle

\begin{abstract}
Credal Variable Elimination (CVE) is an exact algorithm for computing lower and upper
bounds on marginal probabilities in credal networks. It extends the classical variable
elimination algorithm for Bayesian networks to operate on \emph{potentials}---finite
sets of real-valued functions representing the extreme points of local credal sets.
By combining, marginalizing, and pruning these potentials along an elimination ordering,
CVE produces tight probability intervals for any query variable given evidence.
\end{abstract}

\section{Preliminaries}

\subsection{Credal Networks}

A \textbf{credal network} is a pair $\langle G, \mathcal{K} \rangle$ where $G$ is a
directed acyclic graph (DAG) over variables $\mathbf{X} = \{X_1, \dots, X_n\}$ with
finite domains $\mathcal{D} = \{D_1, \dots, D_n\}$, and $\mathcal{K}$ is a collection
of conditional credal sets $K(X_i \mid \mathrm{Pa}(X_i) = \pi)$, one for each variable
$X_i$ and each configuration $\pi$ of its parents $\mathrm{Pa}(X_i)$ in $G$.

Each credal set $K(X_i \mid \pi)$ is a closed convex set of probability distributions
over $X_i$, specified by its \emph{extreme points} (vertices):
\[
  \mathrm{ext}\, K(X_i \mid \pi) = \{P_1(X_i \mid \pi), \dots, P_{v}(X_i \mid \pi)\}.
\]

\subsection{Strong Extension}

The \textbf{strong extension} of a separately specified credal network is the convex
hull of all joint distributions formed by combining one extreme point per local credal
set:
\begin{equation}\label{eq:strong-ext}
  \mathrm{co}\,K(\mathbf{X}) = \mathrm{CH}\Bigl\{
    P(\mathbf{x}) = \prod_{i=1}^{n} P(x_i \mid \pi_i) :
    P(X_i \mid \pi_i) \in \mathrm{ext}\,K(X_i \mid \pi_i)
  \Bigr\}.
\end{equation}

\begin{theorem}[Mau\'a \& Cozman, 2020]
Any combination of extreme points from the local credal sets is an extreme point of the
strong extension. Thus the strong extension can be built directly from the local extreme
points.
\end{theorem}

\subsection{Marginal Inference}

Given a query variable $Z$ and evidence $\mathbf{Y} = \mathbf{y}$, the goal is to compute:
\begin{equation}\label{eq:marginal}
  \underline{P}(Z = z \mid \mathbf{Y} = \mathbf{y}) = \min_{P \in \mathrm{ext}\,K(\mathbf{X})}
  \frac{\sum_{\mathbf{x}'} \prod_i P(x_i \mid \pi_i)}
       {\sum_{z'} \sum_{\mathbf{x}'} \prod_i P(x_i \mid \pi_i)},
\end{equation}
where the sums in numerator and denominator are over the marginalized variables
$\mathbf{X}' = \mathbf{X} \setminus (\{Z\} \cup \mathbf{Y})$, and similarly for the
upper probability $\overline{P}$ with $\max$.

\section{Potentials}

\begin{definition}[Potential]
A \textbf{potential} $\phi(\mathbf{Y})$ over variables $\mathbf{Y}$ is a finite set
of non-negative real-valued functions $\{f_1, f_2, \dots, f_m\}$ on the joint domain
of $\mathbf{Y}$. Each function $f_k$ maps configurations of $\mathbf{Y}$ to
$\mathbb{R}_{\geq 0}$.
\end{definition}

Three operations on potentials are central to the algorithm:

\paragraph{Product.}
The product of $\phi(\mathbf{Y})$ and $\psi(\mathbf{Z})$ is:
\[
  \phi \cdot \psi = \{f \cdot g : f \in \phi, \; g \in \psi\},
\]
where $(f \cdot g)(\mathbf{y}, \mathbf{z}) = f(\mathbf{y}) \cdot g(\mathbf{z})$
over the joint domain $\mathbf{Y} \cup \mathbf{Z}$ (with broadcasting for
non-shared variables).

\paragraph{Sum-Marginal.}
The sum-marginal of $\phi(\mathbf{Y})$ with respect to $X \subseteq \mathbf{Y}$ is:
\[
  \sum_{X} \phi(\mathbf{Y}) = \Bigl\{ \sum_{x} f(\mathbf{y}) : f \in \phi \Bigr\},
\]
producing functions over $\mathbf{Y} \setminus \{X\}$.

\paragraph{Pruning.}
The pruning operation removes dominated functions. We say $g$ \emph{dominates}
$f$ when $g$ is pointwise at least $f$ and strictly greater somewhere:
\begin{definition}[Dominance and pruning]\label{def:prune}
\[
  \prune(\phi) = \bigl\{f \in \phi : \nexists\, g \in \phi \setminus \{f\}
  \text{ s.t.\ } g(\by) \geq f(\by)\;\forall\, \by
  \;\text{ and }\; g(\by') > f(\by') \text{ for some } \by' \bigr\}.
\]
\end{definition}
Because a dominated $f$ can only ever attain the min (resp.\ a dominating $g$
the max) of a quantity that is monotone in the function values, discarding $f$
reduces the cardinality of the potential without affecting the min/max bounds.

\begin{remark}[What the implementation does]\label{rem:prune}
The strictness requirement ``$g > f$ somewhere'' matters, and three properties of
\texttt{Potential.prune} (\texttt{potentials.py:124--148}) are worth recording.
(i)~Two \emph{identical} functions never dominate each other, so $\prune$ is not
a deduplicator: exact duplicates both survive. (ii)~The test is a single greedy
double loop that skips functions already marked dominated, both as candidates and
as dominators, so the surviving set depends on the enumeration order (it is
Pareto-correct, but not a canonical Pareto frontier). (iii)~The comparison is
made on raw floating-point values with \emph{no} tolerance, so numerically
near-identical vertices are retained as distinct functions.
\end{remark}

\section{Algorithm}
\label{sec:algorithm}

CVE answers a query on one variable $Z$ by a single bucket-elimination pass in
which $Z$ is eliminated \emph{last}. Algorithm~\ref{alg:cve} states the pass in
full; the paragraphs that follow it explain the three steps that carry the credal
content.

One structural fact must be stated first, because it shapes the input and the
output. The variables of the credal network are \emph{chain-graph node families},
not the individual atoms of the underlying LCN: a variable $X$ carries a set of
atoms $\mathrm{atoms}(X)$ and has domain size $2^{|\mathrm{atoms}(X)|}$ (the
truth table over its atoms, most-significant bit first). A variable with a single
atom is called \emph{singleton}; one with several is \emph{compound}. The pass
therefore returns bounds on the \emph{states} of $Z$; when $Z$ is singleton its
state-1 bound is the atom marginal directly, and when $Z$ is compound a further
projection is needed to report per-atom marginals
(Section~\ref{sec:projection}).

\begin{algorithm}[H]
\caption{Credal Variable Elimination (CVE)}
\label{alg:cve}
\begin{algorithmic}[1]
\Require Extreme points $\ext\,\cK(X_i \mid \pa(X_i))$ for each variable $X_i$;
         query variable $Z$; evidence $\bY = \by$
\Ensure Lower and upper bounds $[\underline{P}(Z{=}z \mid \by),\;
        \overline{P}(Z{=}z \mid \by)]$ for each state $z$ of $Z$
\Statex
\State $\Gamma \gets \emptyset$
  \Comment{the set of potentials}
\For{each variable $X_i \in \bX$}
  \State $\phi_{X_i} \gets \bigl\{ \text{one chosen }
         P(X_i \mid \pi) \in \ext\,\cK(X_i \mid \pi)
         \text{ per parent config } \pi \bigr\}$
  \State $\Gamma \gets \Gamma \cup \{\phi_{X_i}\}$
    \Comment{$|\phi_{X_i}| = \prod_{\pi} |\ext\,\cK(X_i \mid \pi)|$ CPT vertices}
\EndFor
\For{each evidence variable $Y_j \in \bY$}
  \State $\Gamma \gets \Gamma \cup \bigl\{ \{\delta_{y_j}\} \bigr\}$
    \Comment{indicator potential; the CPTs are not sliced}
\EndFor
\Statex
\State $\bo \gets (X_{\sigma(1)}, \dots, X_{\sigma(n-1)})$, an elimination
       ordering over $\bX \setminus \{Z\}$
  \Comment{\emph{excludes} $Z$}
\For{each $X_i$ in ordering $\bo$}
  \State $\Gamma_{X_i} \gets \{\phi \in \Gamma : X_i \in \scope(\phi)\}$
  \State $\Gamma \gets \Gamma \setminus \Gamma_{X_i}$
  \If{$\Gamma_{X_i} = \emptyset$}
    \State \textbf{continue}
  \EndIf
  \State $\psi \gets \prod \{\phi \in \Gamma_{X_i}\}$
    \Comment{combine all potentials in the bucket}
  \State $\psi' \gets \sum_{X_i} \psi$
    \Comment{marginalize out $X_i$; cardinality unchanged}
  \State $\psi'' \gets \prune(\psi')$
    \Comment{remove dominated functions (Def.~\ref{def:prune})}
  \State $\Gamma \gets \Gamma \cup \{\psi''\}$
\EndFor
\Statex
\State $\phi_{\mathrm{final}} \gets \prod \{\phi \in \Gamma\}$
  \Comment{remaining potentials, scope $\{Z\}$}
\State $\underline{P}(z) \gets 1,\; \overline{P}(z) \gets 0$ for each state $z$
\For{each $f \in \phi_{\mathrm{final}}$ with $\sum_{z'} f(z') > 0$}
  \State $\mathbf{p} \gets f / \sum_{z'} f(z')$
    \Comment{normalize \emph{this} vertex}
  \For{each state $z$}
    \State $\underline{P}(z) \gets \min(\underline{P}(z),\; \mathbf{p}(z))$
    \State $\overline{P}(z) \gets \max(\overline{P}(z),\; \mathbf{p}(z))$
  \EndFor
\EndFor
\State \Return $[\underline{P}(z),\; \overline{P}(z)]$ for each state $z$
\end{algorithmic}
\end{algorithm}

\paragraph{Initial potentials are explicit vertex enumerations.}
Line 3 is where the credal information enters. A family $X_i$ with parent
configurations $\pi_1, \dots, \pi_r$ contributes one function per element of
$\ext\,\cK(X_i \mid \pi_1) \times \cdots \times \ext\,\cK(X_i \mid \pi_r)$: each
function is a complete CPT obtained by choosing one extreme point in every
column, so $|\phi_{X_i}|$ is the \emph{product} of the per-column vertex counts.
The potential is therefore an explicit enumeration of that family's contribution
to the strong extension, which is what makes CVE a strong-extension method.

\paragraph{The query is eliminated last.}\label{sec:querylast}
The ordering on line 9 omits $Z$ entirely, so no step ever marginalizes the query
and its bucket survives to the end. This is the decisive scheduling choice: it is
what makes the answer exact with respect to the strong extension, and it is
exactly what a single shared two-pass ordering (as in CCTE) gives up. Answering
all marginals means repeating the pass once per target, each with its own
ordering, since nothing is shared between targets.

\paragraph{Normalization \emph{is} the evidence conditioning.}
There is no separate division by $P(\by)$. With the indicator potentials of
line 7 in the product, a surviving function of $\phi_{\mathrm{final}}$ satisfies
$f(z) = P(Z{=}z, \by)$ for one selection of extreme points, so
$f(z) / \sum_{z'} f(z') = P(Z{=}z \mid \by)$ for that same selection. Because
each vertex is normalized \emph{before} the min/max is taken, the result is the
correct set of per-vertex conditionals --- not the invalid ratio
$\min P(Z,\by) / \max P(\by)$ of a separately minimized numerator and maximized
denominator. Vertices of zero mass are skipped.

\subsection{Elimination Ordering}

Two heuristics are supported; both exclude the query $Z$, which is what keeps
$Z$'s bucket alive.

\begin{itemize}
  \item \textbf{Topological order}: the DAG's \emph{forward} topological order,
        restricted to the non-query, non-evidence variables, with the evidence
        variables appended at the end. Evidence variables are still eliminated
        --- they are pinned by their indicator potentials rather than removed
        from the problem.
  \item \textbf{Min-fill}: the greedy interaction-graph heuristic, eliminating at
        each step the variable that adds the fewest fill-in edges, with an
        alphabetical tiebreak, and with $Z$ held out of the ordering (it still
        counts as a neighbour when fill is scored). This aims at a small induced
        width $w^*$.
\end{itemize}

\subsection{Projecting a Compound Query onto its Atoms}
\label{sec:projection}

When $Z$ is singleton, its atom marginal is the state-1 entry of the bounds
returned by Algorithm~\ref{alg:cve}. When $Z$ is compound, with $m$ atoms and
$k = 2^m$ states, each atom $a$ is recovered from the state bounds
$[\underline{P}(z), \overline{P}(z)]$ by two linear programs over the states
$S_a$ in which $a$ is true:
\[
  \underline{P}(a) = \min \Bigl\{ \sum_{z \in S_a} p_z \;:\;
    \sum_{z} p_z = 1, \;\;
    \underline{P}(z) \le p_z \le \overline{P}(z) \Bigr\},
\]
and symmetrically $\overline{P}(a)$ with $\max$. The feasible region is the box
$\prod_z [\underline{P}(z), \overline{P}(z)]$ intersected with the probability
simplex, which is an \emph{outer relaxation} of the true state polytope of $Z$.
Atom bounds on a compound query can therefore be looser than the state bounds
CVE actually computed; the running example in Section~\ref{sec:example} shows
this at its most extreme.

\section{Running Example}
\label{sec:example}

\begin{example}[Alarm Network]
The LCN \texttt{examples/alarm.lcn} over atoms $A, B, C, D, E$ factorizes into
four credal-network nodes: the singletons $B$ (Burglary), $E$ (Earthquake),
$A$ (Alarm), and the \emph{compound} node $C\text{-}D$ carrying the two atoms
$C, D$ (they are merged because the sentence
$0.5 \le P(\neg(C \oplus D) \mid A) \le 0.7$ constrains them jointly). The
compound node has $2^2 = 4$ states, ordered $CD \in \{00, 01, 10, 11\}$
most-significant bit first. The DAG is
\[
  B \to A, \quad E \to A, \quad A \to C\text{-}D .
\]
The local credal sets (extreme points from LRS vertex enumeration) are:

\begin{center}
\begin{tabular}{llr@{\hskip 1em}l}
\toprule
\textbf{Node} & \textbf{Parent Config} & \textbf{\#} & \textbf{Vertices} \\
\midrule
$B$ & --- & 2 & $[0.8, 0.2]$, $[0.9, 0.1]$ \\
$E$ & --- & 2 & $[0.9, 0.1]$, $[0.95, 0.05]$ \\
$A$ & $B{=}0, E{=}0$ & 2 & $[0.956, 0.044]$, $[1.0, 0.0]$ \\
$A$ & $B{=}1, E{=}0$ & 2 & $[0.0, 1.0]$, $[1.0, 0.0]$ \\
$A$ & $B{=}0, E{=}1$ & 2 & $[0.0, 1.0]$, $[1.0, 0.0]$ \\
$A$ & $B{=}1, E{=}1$ & 2 & $[0.0, 1.0]$, $[1.0, 0.0]$ \\
$C\text{-}D$ & $A{=}0$ & 4 & $[0,0,0,1]$, $[0,0,1,0]$, $[0,1,0,0]$, $[1,0,0,0]$ \\
$C\text{-}D$ & $A{=}1$ & 11 & $[0,0,0.3,0.7]$, $[0,0.3,0,0.7]$, \dots, $[0.7,0,0,0.3]$ \\
\bottomrule
\end{tabular}
\end{center}

\paragraph{Initial potentials.}
Line 3 of Algorithm~\ref{alg:cve} takes one extreme point per parent
configuration, so the potential sizes are the \emph{products} of the counts
above:
\[
  |\phi_B| = 2, \quad |\phi_E| = 2, \quad
  |\phi_A| = 2^4 = 16, \quad
  |\phi_{C\text{-}D}| = 4 \times 11 = 44 ,
\]
for $2 + 2 + 16 + 44 = 64$ functions in total across the four potentials.

\paragraph{Query: $P(B{=}1)$ (no evidence).}
The forward topological order with the target removed is
$\bo = (E, A, C\text{-}D)$; note $B$ never appears in it.

\textbf{Step 1: eliminate $E$.}
The bucket is $\phi_E$ (scope $\{E\}$, 2 functions) and $\phi_A$ (scope
$\{A,B,E\}$, 16 functions). Their product has $2 \times 16 = 32$ functions over
$\{A,B,E\}$; summing out $E$ leaves 32 functions over $\{A,B\}$, and pruning
removes none of them.

\textbf{Step 2: eliminate $A$.}
The bucket is the new potential over $\{A,B\}$ (32 functions) together with
$\phi_{C\text{-}D}$ (scope $\{C\text{-}D, A\}$, 44 functions), giving
$32 \times 44 = 1408$ functions over $\{C\text{-}D, A, B\}$. Summing out $A$
leaves 1408 functions over $\{C\text{-}D, B\}$; again none are dominated.

\textbf{Step 3: eliminate $C\text{-}D$.}
Summing out $C\text{-}D$ leaves 1408 functions over $\{B\}$, and here pruning
bites: only $328$ survive.

\textbf{Result.} Each surviving $f$ over $\{B\}$ is normalized,
$p(B{=}1) = f(1)/(f(0)+f(1))$, and the componentwise min/max over the $328$
normalized vectors gives
\[
  \underline{P}(B{=}1) = 0.10, \qquad \overline{P}(B{=}1) = 0.20 ,
\]
recovering exactly the interval asserted by the sentence
$0.1 \le P(B) \le 0.2$, as it must since $B$ is a root.

\paragraph{The compound node illustrates the projection gap.}
Running the driver over all four nodes returns $P(A{=}1) \in [0, 0.312]$ and
$P(E{=}1) \in [0.05, 0.10]$, but the state bounds for $C\text{-}D$ come out
vacuous, $[\underline{p}, \overline{p}] = [\mathbf{0}, \mathbf{1}]$ on all four
states. The per-atom linear programs of Section~\ref{sec:projection} can then
only return $P(C{=}1), P(D{=}1) \in [0,1]$: with every
state interval equal to $[0,1]$, the box--simplex relaxation carries no
information, even though the joint constraint on $C \oplus D$ does. This is the
outer relaxation of the compound-node projection at its worst.
\end{example}

% =====================================================================
% --- self-contained; safe to lift into a paper submission ---
% Everything below is written to stand alone: it repeats the definitions it
% needs and refers to no other part of this document. Wording is kept aligned
% with docs/neurips_2026.tex, Section "Credal Variable Elimination", so the two
% cannot drift apart. Prose only -- no pseudocode, no floats; requires only the
% \newcommand macros in the preamble (or docs/macros.tex).
% =====================================================================
\section{Succinct Description (paper-ready)}
\label{sec:succinct}

% When lifting into a paper, replace the author-year mention below with
% \cite{dechter1999bucket}.
Credal Variable Elimination (CVE) extends classical bucket elimination
(Dechter, 1999) to credal networks by replacing every factor with a
\emph{potential}: a finite \emph{set} of non-negative functions over a common
scope, each member being one extreme point of the corresponding local credal set.
Potentials support three operations. Their \emph{product} is the elementwise
product of every pair, $\phi \cdot \psi = \{f \cdot g : f \in \phi,\,
g \in \psi\}$, so cardinalities multiply; the \emph{sum-marginal}
$\sum_X \phi = \{\sum_x f : f \in \phi\}$ sums each member separately and leaves
the cardinality unchanged; and \emph{pruning} discards any $f$ for which some
$g \in \phi$ satisfies $g \geq f$ pointwise with $g > f$ somewhere, which shrinks
the set without affecting its extremes, since a dominated function attains
neither. To answer a query on $Z$ under evidence $\bY = \by$, CVE builds one
potential per family --- enumerating the conditional probability tables obtained
by choosing one extreme point in each parent configuration, so the potential
explicitly lists that family's contribution to the strong extension --- adds a
single-function indicator potential for each observation, and then eliminates the
variables one at a time along an ordering that \emph{excludes} $Z$: at each step
it multiplies together all potentials whose scope contains the current variable,
sums that variable out, and prunes the result.

Because $Z$ is never eliminated, its bucket survives to the end, and the final
product is a potential over $Z$ alone. Each of its members equals
$P(Z, \by)$ under one selection of extreme points, so dividing that member by its
own sum yields $P(Z \mid \by)$ for the same selection; taking the componentwise
minimum and maximum \emph{after} normalizing each member therefore gives bounds
that are exact with respect to the strong extension --- and not the invalid ratio
of a separately minimized numerator and a separately maximized denominator.
Never eliminating the query is what buys this exactness, and it is precisely what
a single shared two-pass ordering gives up; the price is that bounds for all
variables require one pass per target, with nothing reused between them. The cost
of a pass is $O(n \cdot C \cdot k^{w^*})$ for $n$ variables of domain size at
most $k$, an ordering of induced width $w^*$, and $C$ bounding the cardinality of
the intermediate potentials. Pruning is what keeps $C$ in check: it can still
grow exponentially in the number of extreme points in the worst case, but remains
manageable in practice on networks of moderate induced width.
% =====================================================================
% --- end of self-contained block ---
% =====================================================================

\section{Complexity}

\begin{theorem}
The time and space complexity of CVE is $O(n \cdot C \cdot k^{w^*})$, where $n$ is
the number of variables, $k$ bounds the domain sizes, $w^*$ is the induced width of the
elimination ordering, and $C$ bounds the cardinality of the intermediate potentials.
\end{theorem}

In the worst case, $C$ is exponential in the number of extreme points. However, pruning
keeps $C$ manageable in practice for networks of moderate treewidth.

\section{References}

\begin{enumerate}
\item D.D.\ Mau\'a, F.G.\ Cozman. Thirty years of credal networks: Specification,
      algorithms and complexity. \emph{International Journal of Approximate Reasoning},
      126:133--157, 2020.
\item R.\ Marinescu, R.\ Dechter, A.\ Ihler. Credal Marginal MAP. \emph{NeurIPS}, 2023.
\item F.G.\ Cozman. Credal networks. \emph{Artificial Intelligence}, 120(2):199--233, 2000.
\item R.\ Dechter. Bucket elimination: A unifying framework for reasoning.
      \emph{Artificial Intelligence}, 113(1--2):41--85, 1999.
\item D.D.\ Mau\'a, C.P.\ de Campos, A.\ Benavoli, A.\ Antonucci. Probabilistic
      inference in credal networks: New complexity results. \emph{Journal of
      Artificial Intelligence Research}, 50:603--637, 2014.
\end{enumerate}

\end{document}
